\documentclass[11pt]{article}
\usepackage[a4paper,margin=27mm]{geometry}
\usepackage[T1]{fontenc}
\usepackage{lmodern,amsmath,amssymb,amsthm,booktabs,array,longtable}
\usepackage[colorlinks=true,linkcolor=blue,citecolor=blue,urlcolor=blue]{hyperref}
\newtheorem{theorem}{Theorem}
\newtheorem{lemma}[theorem]{Lemma}
\newtheorem{corollary}[theorem]{Corollary}
\newcommand{\coeff}[2]{[#1]#2}
\newcommand{\Q}{\mathbb{Q}}
\title{Mathar's path recurrences and alternative proofs\\of Prodinger's enumeration}
\author{Haobo Ma\thanks{ChronoAI Pte Ltd, Singapore.}
\and Wenlin Zhang\thanks{National University of Singapore; ChronoAI Pte Ltd, Singapore.}}
\date{17 November 2026}
\begin{document}
\maketitle
\begin{center}\small Licensed under Creative Commons Attribution 4.0 International (CC BY 4.0).\end{center}
\begin{abstract}
Using Prodinger's enumeration (arXiv:1912.06930v3, 2020, Theorem 1 and
Section 3), we give alternative counting proofs and prove Mathar's eight-term
recurrence for floor $-2$ paths (OEIS A379464, revision 20, second Formula,
29 January 2025, $n\geq7$) and his six-term recurrence for floor $-1$
(A379463, revision 16, second Formula, same date, $n\geq5$).
An explicit differential-operator certificate transports the known Raney counts
to the exact recurrences. We include every fixed lower floor by a finite
polynomial correction, prove the recorded asymptotics, and correct the
indexing and tree wording in A127902 and companion formulas in A069271.
\end{abstract}

\section{The paths and the recurrence conjectures}
For $r\in\{1,2,3\}$ and $n\geq0$, let $a_r(n)$ count words of length $n$ in
$\{1,0,-3\}$ with total sum zero and every prefix sum at least $1-r$.
The empty word is counted once. Thus $a_3(n)$ is A379464, $a_2(n)$ is
A379463, and $a_1(n)$ is A127902 at index $n+1$ \cite{oeis3,oeis2,oeis1}.
These are respectively the lattice paths with steps $(1,1),(1,0),(1,-3)$
whose floors are $-2,-1,0$.

The two source lines are headed ``Conjecture D-finite with recurrence'' and
signed R.\ J.\ Mathar on 29 January 2025. Their locators are A379464
revision 20, second Formula, and A379463 revision 16, second Formula
(split across two text-export lines). The coefficients are reproduced in
Appendix~\ref{sec:coefficients}. Scheuerle recorded the two binomial programs
on 7 January 2025; these formulas, rather than newly discovered terms, are
the starting point.

The zero-free counts are known: the Raney formula follows already from
Krattenthaler's Theorem 10.4.5 and its cycle-lemma proof
\cite{krattenthaler}. Prodinger's Theorem 1 credits Selkirk's
2019 thesis and proves the lower-floor enumeration; his Section 3 treats
arbitrary floors \cite{prodinger}. Drube later gives a general starting- and
ending-height formula \cite[Theorem 2.3 and Corollary 2.4]{drube}.
We reprove the counting statement, insert horizontal steps, and supply the
explicit certificates for Mathar's recurrences. The all-floor formula below
is a specialization of this earlier enumeration, not a new enumeration theorem.
The corrections concern the literal formulas and conventions in the cited
OEIS revisions.

Asinowski, Hackl and Selkirk give the same prepend-and-rotate counting
proof, including its periodic-rotation multiplicities
\cite[Proposition 3, equation (2)]{ahs}; DeJager, Seidl, Naquin and Drube
obtain the corresponding generating-function powers through Riordan arrays
\cite[Corollary 11]{dejager}. Barry explicitly derives the zero-floor
binomial sum in Example 16 of his article available online in November
2015 \cite{barry}. The ordered-tree correspondence and its asymptotic
specialization also follow from Banderier and Flajolet's earlier
tree-code framework \cite[Examples 3 and 6]{bf}.
Differential equations followed by coefficient extraction are standard
recurrence tools; Mathar treats this general approach in \cite{mathar2021}.
Here the contribution is the explicit operator identities for the stated
path counts and source recurrences.

Put $F_r(x)=\sum_{n\geq0}a_r(n)x^n$. All generating-function manipulations
before Section~\ref{sec:asymptotic} take place in formal power series over $\Q$.

\begin{theorem}\label{thm:main}
For $r=1,2,3$ and $n\geq0$,
\begin{equation}\label{eq:count}
 a_r(n)=\sum_{k=0}^{\lfloor n/4\rfloor}\binom n{4k}
             \frac{r}{4k+r}\binom{4k+r}{k}.
\end{equation}
Let $C(t)$ be the unique formal series with $C(0)=1$ satisfying
$C(t)=1+tC(t)^4$. Then
\begin{equation}\label{eq:gf}
 F_r(x)=\frac1{1-x}C\!\left(\frac{x^4}{(1-x)^4}\right)^r.
\end{equation}
The coefficients $p_{r,j}(n)$ in Appendix~\ref{sec:coefficients} satisfy
\begin{equation}\label{eq:recurrences}
 \sum_{j=0}^{d_r}p_{r,j}(n)a_r(n-j)=0
 \quad(n\geq d_r),\qquad (d_1,d_2,d_3)=(4,5,7).
\end{equation}
For $r=3$ and $r=2$ these are exactly Mathar's eight-term and six-term
conjectured recurrences, respectively. For $r=1$ this is Kotesovec's
five-term recurrence after the change of index $n\mapsto n+1$.
\end{theorem}

For example, $a_3(0),\ldots,a_3(7)=1,1,1,1,4,16,46,106$.
At length four the admissible words are $0000$, $1,-3,1,1$,
$1,1,-3,1$, and $1,1,1,-3$. Since $p_{3,0}(n)>0$ for $n\geq7$,
the first seven values and the eight-term recurrence determine $a_3$ uniquely.
Likewise $p_{2,0}(n)>0$ for $n\geq5$, so the first five values determine
$a_2$ from its recurrence.

\section{Counting by cyclic shifts}
We give the version of the cycle lemma needed to identify the literal path
count. Cuts are positions, so periodic words cause no division ambiguity.

\begin{lemma}\label{lem:cycle}
Let $w_1,\ldots,w_N$ be integers at most one, with sum $r>0$.
Exactly $r$ of the $N$ cyclic cuts give a word whose nonempty prefix sums
are all strictly positive, counting cuts with their positions.
\end{lemma}
\begin{proof}
Write $S_0=0$, $S_i=w_1+\cdots+w_i$, and $m=\min_{0\leq i\leq N}S_i$.
A cut after position $j\in\{0,\ldots,N-1\}$ is good precisely when
\begin{equation}\label{eq:goodcut}
 S_i>S_j\ (j<i\leq N),\qquad S_i+r>S_j\ (0\leq i\leq j).
\end{equation}
For each integer $c\in\{m,\ldots,m+r-1\}$, the path of partial sums
visits $c$: from an occurrence of its minimum to $S_N=r$, upward jumps
have size at most one. Let $j_c$ be its last visit to $c$.
As $c<r$, this visit precedes $N$. After $j_c$ every partial sum exceeds
$c$, since a later value below $c$ would require another visit on the way
to $r$. Before $j_c$ every partial sum is at least $m$, and
$m+r>c$. Hence \eqref{eq:goodcut} holds at $j_c$.

Conversely, if a cut is good, a minimum cannot occur after it. Applying
the second inequality in \eqref{eq:goodcut} at a minimum gives
$m\leq S_j<m+r$; the first inequality makes $j$ the last visit to its
level. Thus the good cuts are exactly the $r$ positions $j_c$.
\end{proof}

Erase the zero steps from an admissible word. If $k$ steps are $-3$,
there are $3k$ steps $1$, so the remaining word has length $4k$.
Prepending $r$ steps $1$ converts it into a word of length $4k+r$,
sum $r$, and strictly positive nonempty prefix sums.
Conversely, a positive word in $\{1,-3\}$ starts with at least $r$
steps $1$ when $r\leq3$: a $-3$ among its first three positions would
make a prefix sum nonpositive. Removing the first $r$ steps gives
exactly a zero-sum word with floor $1-r$.

There are $\binom{4k+r}{k}$ unrestricted words of length $4k+r$ and
sum $r$. Count pairs consisting of such a word and a good cyclic cut.
Lemma~\ref{lem:cycle} gives $r\binom{4k+r}{k}$ pairs. Each positive word
has exactly $4k+r$ inverse pairs, one for each cut position, even if
some of the rotated words coincide. Consequently the number of zero-free
words is
\begin{equation}\label{eq:b}
 b_r(k)=\frac{r}{4k+r}\binom{4k+r}{k}
       =\frac{r(4k+r-1)!}{k!(3k+r)!}.
\end{equation}
Choose the $4k$ nonzero positions among $n$ positions. Reinserting zeros
does not change any attained height, which proves \eqref{eq:count}.

For completeness, the formal Lagrange inversion formula
\cite[Section I.5, formula (69)]{ac}, applied to $W=t(1+W)^4$,
gives, for $k\geq1$,
\[
 [t^k](1+W)^r=\frac r k [u^{k-1}](1+u)^{4k+r-1}
             =\frac r{4k+r}\binom{4k+r}{k}.
\]
The constant coefficient is one. Hence $\sum_{k\geq0}b_r(k)t^k=C(t)^r$.
Finally $\sum_{n\geq m}\binom nm x^n=x^m/(1-x)^{m+1}$ proves
\eqref{eq:gf}. This finishes the counting part of Theorem~\ref{thm:main}.

\begin{corollary}[All lower floors; Prodinger's formula with horizontal steps]
\label{cor:allfloors}
For an integer $h\geq0$, let $f_h(n)$ count zero-sum words of length $n$
in $\{1,0,-3\}$ with every prefix sum at least $-h$. Define
\[
 D_h(t)=\sum_{i=0}^{\lfloor h/4\rfloor}(-1)^i\binom{h-3i}{i}t^i,
 \qquad t(x)=\frac{x^4}{(1-x)^4}.
\]
Then
\begin{equation}\label{eq:allfloors}
 \sum_{n\geq0} f_h(n)x^n=
 \frac{D_h(t(x))C(t(x))^{h+1}}{1-x}.
\end{equation}
In particular, the Raney sum \eqref{eq:count} holds also for $r=4$
(floor $-3$). For floor $-4$ the zero-free generating function is
$(1-t)C(t)^5$, not $C(t)^5$.
\end{corollary}
\begin{proof}
This is the horizontal-step form of Prodinger's Section 3 formula
\cite{prodinger}, also Drube's equation (4) with $\alpha=\beta=h$
\cite{drube}. Here is a formal derivation. Let $G_{\alpha,\beta}(t)$
count words in $\{1,-3\}$ from height $\alpha$ to height $\beta$, staying
nonnegative, with $t$ marking down-steps; take $G_{\alpha,\beta}=0$ for
$\alpha<0$. Last passages through levels $0,\ldots,\beta-1$ decompose
a path from $0$ to $\beta$ into $\beta+1$ nonnegative excursions and
$\beta$ separating up-steps. Each excursion has generating function $C$:
a first-return decomposition gives $C=1+tC^4$. Thus
$G_{0,\beta}=C^{\beta+1}$.

Decomposing the first step from height $\alpha-1$, for $\alpha\geq1$,
gives
\[
 G_{\alpha,\beta}=G_{\alpha-1,\beta}-tG_{\alpha-4,\beta}
                   -\boldsymbol{1}_{\alpha-1=\beta}.
\]
For $H_\beta(q,t)=\sum_{\alpha\geq0}G_{\alpha,\beta}(t)q^\alpha$
this says
\[
 H_\beta(q,t)=\frac{C(t)^{\beta+1}-q^{\beta+1}}{1-q+tq^4}.
\]
The coefficient of $q^h$ in $1/(1-q+tq^4)$ is $D_h(t)$, by expansion
as $\sum_{j\geq0}(q-tq^4)^j$. Set $\alpha=\beta=h$: the numerator
term $q^{h+1}$ contributes nothing, so the zero-free generating function
is $D_h C^{h+1}$. Inserting zero steps gives \eqref{eq:allfloors}.
For $0\leq h\leq3$, $D_h=1$; for $h=4$, $D_4=1-t$.
The naive extension to $r=5$ is not asserted in the posing source and
fails at length four: floor $-4$ permits five words, while that sum gives six.
\end{proof}

\section{From factorial ratios to the exact recurrences}
Let $\theta_t=t\,d/dt$, $\theta=x\,d/dx$, and define
\[
 E=x(1-x)\frac d{dx}-x,\qquad
 t(x)=\frac{x^4}{(1-x)^4},\qquad
 (\mathcal T B)(x)=\frac{B(t(x))}{1-x}.
\]
The chain rule gives
\begin{equation}\label{eq:conjugacy}
 E\mathcal T B=4\mathcal T(\theta_t B),
\end{equation}
because $x(1-x)(1/(1-x))'-x/(1-x)=0$ and
$x(1-x)t'(x)=4t(x)$.

The factorial expression \eqref{eq:b} yields
\[
 \frac{b_r(k+1)}{b_r(k)}=
 \frac{\prod_{h=0}^{3}(4k+r+h)}{(k+1)\prod_{h=1}^{3}(3k+r+h)}.
\]
Cancelling factors gives $P_r(k+1)b_r(k+1)=Q_r(k)b_r(k)$, where
\begin{center}\small
\begin{tabular}{c ll}
\toprule
$r$&$P_r(z)$&$Q_r(z)$\\\midrule
1&$3z(3z-1)(3z+1)$&$8(2z+1)(4z+1)(4z+3)$\\
2&$3z(3z+1)(3z+2)$&$8(2z+1)(4z+3)(4z+5)$\\
3&$3(z+1)(3z+1)(3z+2)$&$8(2z+3)(4z+3)(4z+5)$\\
\bottomrule
\end{tabular}
\end{center}
As $b_r(0)=1$, coefficient comparison gives
\begin{equation}\label{eq:base}
 (P_r(\theta_t)-tQ_r(\theta_t))C(t)^r=P_r(0).
\end{equation}
Define the operators
\[
 L_r=(1-x)^4U_r(E)-256x^4V_r(E),
 \quad U_r(z)=64P_r(z/4),\quad V_r(z)=Q_r(z/4)/4.
\]
Thus
\begin{center}\small
\begin{tabular}{c ll}
\toprule
$r$&$U_r(z)$&$V_r(z)$\\\midrule
1&$27z^3-48z$&$z^3+6z^2+11z+6$\\
2&$27z^3+108z^2+96z$&$z^3+10z^2+31z+30$\\
3&$27z^3+216z^2+528z+384$&$z^3+14z^2+63z+90$\\
\bottomrule
\end{tabular}
\end{center}
Equations \eqref{eq:gf}--\eqref{eq:base} imply
\begin{equation}\label{eq:forcing}
 L_1F_1=L_2F_2=0,\qquad L_3F_3=384(1-x)^3.
\end{equation}
Indeed, $L_r\mathcal T B=64(1-x)^4\mathcal T
((P_r(\theta_t)-tQ_r(\theta_t))B)$, and
$P_1(0)=P_2(0)=0$, $P_3(0)=6$.

Here is an explicit coefficient certificate for the final step. Set
\[
 R_r=\sum_{j=0}^{d_r}x^j p_{r,j}(\theta+j),
\]
with the polynomials in Appendix~\ref{sec:coefficients}. Then
\begin{equation}\label{eq:operators}
 R_1=(1-x)^{-3}L_1,\qquad
 R_2=(1-x)^{-2}(\theta-3)L_2,\qquad R_3=L_3.
\end{equation}
Operator products act from right to left; the two inverse powers mean
multiplication by formal series with constant coefficient one.
To expand these identities one uses $Dx=xD+1$, where $D=d/dx$.
More explicitly, if $f=\sum f_nx^n$, then for $n\geq3$
\begin{align*}
 [x^n]Ef&=n(f_n-f_{n-1}),\\
 [x^n]E^2f&=n^2f_n-n(2n-1)f_{n-1}+n(n-1)f_{n-2},\\
 [x^n]E^3f&=n^3f_n-n(3n^2-3n+1)f_{n-1}
             +3n(n-1)^2f_{n-2}-n(n-1)(n-2)f_{n-3}.
\end{align*}
Substitution of these three lines into $L_r$, with
$(1-x)^4=1-4x+6x^2-4x^3+x^4$, gives \eqref{eq:operators}.
For $R_2$ the product rule is applied once more to $(\theta-3)L_2$;
multiplying by $(1-x)^2$ clears its displayed denominator.
This supplies every coefficient in the appendix. For instance the extreme
coefficients for $r=3$ are $3(n+4)(3n+4)(3n+8)$ and
$229(n-4)(n-5)(n-6)$.

For $n\geq d_r$, the coefficient of $x^n$ in $R_rF_r$ is exactly the
left side of \eqref{eq:recurrences}. The forcing term in
\eqref{eq:forcing} has degree three. Taking coefficients therefore proves
all three recurrences and completes Theorem~\ref{thm:main}.

\section{The recorded asymptotic formulas}\label{sec:asymptotic}
\begin{corollary}\label{cor:asymptotic}
Put $\lambda=1+4/3^{3/4}$, $\rho=\lambda^{-1}$, and $\tau=4/3$.
For $r=1,2,3$, as $n\to\infty$,
\begin{equation}\label{eq:asymptotic}
 a_r(n)\sim
 \frac{r\tau^{r-1}\sqrt{8/27}}{\sqrt\pi(1-\rho)^{3/2}}
 \frac{\lambda^n}{n^{3/2}}.
\end{equation}
In particular, $a_3(n)\sim
2^{5/2}\lambda^{n+3/2}/(3^{11/8}\sqrt\pi\,n^{3/2})$,
and the constant for $a_2$ is half as large.
\end{corollary}
\begin{proof}
The characteristic equations for $C=1+tC^4$ give
$t_0=27/256$, $C(t_0)=\tau=4/3$.
Writing $C=\tau-\alpha\sqrt{1-t/t_0}+O(1-t/t_0)$ in that equation,
the linear term in $C-\tau$ vanishes and the next terms give
$\alpha^2=\tau^2/6=8/27$. The branch with $C(0)=1$ has
$\alpha>0$. This local expansion converges in powers of
$\sqrt{1-t/t_0}$ by the analytic implicit-function theorem after this
change of variable.

The only nonzero finite critical value of this algebraic equation is $t_0$:
solving $C-1-tC^4=0$ and $1-4tC^3=0$ gives the values above.
Under $t=t(x)$ the four possible critical points are
$x=s\zeta/(1+s\zeta)$, where $s=3^{3/4}/4$ and $\zeta^4=1$.
Only $\zeta=1$ has minimal modulus, namely $\rho$; for all other choices
$|1+s\zeta|<1+s$. The point $x=1$ is also farther from zero.
The branch analytic at zero consequently has $\rho$ as its unique
dominant singularity and is analytic in a dented neighborhood there.
Moreover
\[
 1-\frac{t(x)}{t_0}=\frac4{1-\rho}(1-x/\rho)
                           +O((1-x/\rho)^2).
\]
Inserting the expansion of $C$ into \eqref{eq:gf} shows that the coefficient
of $\sqrt{1-x/\rho}$ in $F_r$ is
\[
 -\frac{2r\tau^{r-1}\sqrt{8/27}}{(1-\rho)^{3/2}}.
\]
The transfer formula
$[x^n](1-x/\rho)^{1/2}\sim-\rho^{-n}/(2\sqrt\pi n^{3/2})$
\cite[Theorems VI.1 and VI.3]{ac} proves \eqref{eq:asymptotic};
the full local expansion justifies transfer of its higher-order remainder.
For $r=3$ simplifying with $1-\rho=(4/3^{3/4})/\lambda$ gives the
stated constant. The factor $r\tau^{r-1}$ gives the ratio $1/2$ for
$r=2$ and $3/16$ for $r=1$ relative to $r=3$.
\end{proof}

For A127902 at its printed index $N$, substitute $n=N-1$ in
\eqref{eq:asymptotic}. Its leading constant becomes
$\sqrt{4+3^{3/4}}\,3^{1/4}/(12\sqrt{\pi/2})$ times
$\lambda^N/N^{3/2}$, exactly the form recorded by Kotesovec on
20 August 2013. The constants for A379463 and A379464 agree with his
15 January 2025 formulas.

\section{Companion formulas and corrections}\label{sec:siblings}
\begin{corollary}[The printed index of A127902]\label{cor:offset}
Write $u_N=\mathrm{A127902}(N)$, $N\geq1$. Then
\begin{equation}\label{eq:offset}
 u_N=\sum_{k=0}^{\lfloor(N-1)/4\rfloor}
       \binom{N-1}{4k}\frac{\binom{4k}{k}}{3k+1},
 \qquad Y(x)=\sum_{N\geq1}u_Nx^N=x(1+Y+Y^4).
\end{equation}
Thus $u_N$ counts ordered rooted trees without vertex labels, with $N$
vertices and outdegrees $0,1,4$.
\end{corollary}
\begin{proof}
The identity $u_N=a_1(N-1)$ gives the sum. Substituting
$C=(1-x)F_1$ into its defining equation gives
$(1-x)F_1=1+x^4F_1^4$; setting $Y=xF_1$ proves the second identity.
Root decomposition proves the tree statement, because a root is a leaf,
has one child, or has four ordered children.
\end{proof}
The sum is Barry's zero-based coefficient formula
\cite[Example 16, pp.\ 361--362]{barry}, shifted to the printed OEIS index;
the tree correspondence is the specialization of
\cite[Example 3, equation (24)]{bf} to outdegrees $0,1,4$.
The first Formula of A127902 revision 38 instead prints the sum with $N$
in place of $N-1$; it gives $2$ at $N=4$, whereas $u_4=1$.
Equation~\eqref{eq:offset} is the corrected formula for the printed offset.
The PARI and Magma programs enumerate path lengths starting at zero, so
their first outputs correspond to $u_1$, not $u_0$.
Devlin's 4 March 2012 comment says ``rooted labeled trees''.
Under distinct vertex labeling this fails on vertex set $\{1,2\}$:
either vertex can be the root, giving two trees, whereas $u_2=1$.
The replacement interpretation is the one in Corollary~\ref{cor:offset}.
Heinz's additional upper bound $y\leq x$ is automatic for the increments
$1,0,-3$ and does not change the floor-zero count.

\begin{corollary}[Zero-free siblings]\label{cor:zerofree}
The coefficients $b_1(k)$ and $b_2(k)$ in \eqref{eq:b} are A002293$(k)$
and A069271$(k)$, respectively. Their generating functions are $C(t)$ and
$A(t)=C(t)^2$. In particular,
\begin{gather}\label{eq:A}
 A=(1+tA^2)^2,\qquad
 3k(3k+1)(3k+2)b_2(k)=8(4k+1)(2k-1)(4k-1)b_2(k-1)
 \quad(k\geq1),\\
 b_1(k)\sim\frac{\sqrt2}{3^{3/2}\sqrt\pi}
        \frac{(256/27)^k}{k^{3/2}},\qquad
 b_2(k)\sim\frac{2^{7/2}}{3^{5/2}\sqrt\pi}
        \frac{(256/27)^k}{k^{3/2}}.\label{eq:zeroasymptotic}
\end{gather}
For every real $m$ and integer $n\geq1$,
\begin{equation}\label{eq:powers}
 [t^n]A(t)^{mn}=2m\binom{(2m+4)n-1}{n-1}.
\end{equation}
Here $A^{mn}$ and generalized binomial coefficients have their formal
binomial-series meanings. In particular,
\begin{equation}\label{eq:inverse}
 [t^n]A(t)^{-2n}=4(-1)^n.
\end{equation}
Together with $A(0)=1$, these conditions uniquely determine $A$.
\end{corollary}
\begin{proof}
The coefficient formulas in the two entries \cite{oeis4,oeis5} are
\eqref{eq:b}. Root decomposition shows that $C$ counts full $4$-ary trees
by internal vertices; $C^2$ counts ordered pairs of them.
Thus $A=C^2=(1+tC^4)^2=(1+tA^2)^2$; the recurrence
is the $r=2$ factorial ratio from Section 3 with $k$ shifted by one.
The square-root expansion in Section~\ref{sec:asymptotic}, now in $t$,
has coefficients $-\sqrt{8/27}$ and $-2(4/3)\sqrt{8/27}$ for $C$ and
$C^2$. Transfer gives \eqref{eq:zeroasymptotic}.
Lagrange inversion with $W=t(1+W)^4$ gives, for any real $s$,
\[
 [t^n](1+W)^s=\frac{s}{n}\binom{4n+s-1}{n-1}.
\]
Take $s=2mn$ to prove \eqref{eq:powers}. At $m=-2$ this is
$-4\binom{-1}{n-1}=4(-1)^n$. If $A=1+\sum_{j\geq1}a_jt^j$,
the coefficient of $t^n$ in $A^{-2n}$ is $-2na_n$ plus a polynomial
in $a_1,\ldots,a_{n-1}$; hence the conditions determine $a_n$ successively.
\end{proof}
A069271 revision 178, Bala's 10 September 2025 block, prints $-4$ in
place of \eqref{eq:inverse}. At $n=2$, $A=1+2t+9t^2+\cdots$ gives
$[t^2]A^{-4}=4$, not $-4$. The following displayed power formula there
is valid with its denominator when $m\ne-2$, not the printed exclusion
$m\ne2$; equation~\eqref{eq:powers} covers the exceptional value too.
Somos's PARI line of 28 March 2012 has an extra factor $2$:
the correct expression is $\binom{4n+2}{n}/(2n+1)$, not twice this
quantity (already $n=0$ gives $1$, not $2$).

Other algebraic forms in A069271 follow directly. With formal square roots
having constant coefficient one,
\begin{equation}\label{eq:otherA}
 2t^2A^3=1-2tA-\sqrt{1-4tA},\qquad
 B(z)=zA(z^2),\quad z=\frac{B}{(1+B^2)^2}.
\end{equation}
For the first equation substitute $t=(C-1)/C^4$ and use
$\sqrt{1-4tC^2}=2/C-1$; the second is \eqref{eq:A}.
If $K(v)=1+vK(v)^2$, then $K(tA)=C$, since
$1+tAC^2=1+tC^4=C$. Therefore $tA$ is the reversion of $v/K(v)^2$,
which is Bala's 8 October 2015 formula. Lagrange inversion applied to
$\log C$ also gives
\[
 [t^n]\frac{tA'}{2A}=\binom{4n-1}{n-1}\quad(n\geq1).
\]
The factorial ratio identifies the ordinary hypergeometric series as
${}_3F_2(\frac12,\frac34,\frac54;\frac43,\frac53;\frac{256t}{27})$;
dividing the coefficients by $n!$ adds the denominator parameter $1$
for the exponential generating function. Finally, substituting
$t=z/(1+z)^4$ into $C=1+tC^4$ gives $C=1+z$ and proves Gutkovskiy's
9 September 2026 identity
$\sum_{n\geq0}b_2(n)z^n/(1+z)^{4n+2}=1$.

The production-matrix statements in A069271 and A002293 have the same
short explanation. Put $f=t/(1-f)^3$ and $g=1/(1-f)=C$.
For the matrix with $M_{i0}=\binom{i+2}{2}$ and
$M_{ij}=\binom{i-j+3}{2}$ for $1\leq j\leq i+1$ (zero otherwise),
the column generating functions of $(M^n)_{0j}$ are $V_j=gf^j$.
Indeed, their defining recurrences are
$V_0=1+t\sum_iV_i\binom{i+2}{2}$ and
$V_j=t\sum_{i\geq j-1}V_i\binom{i-j+3}{2}$ for $j\geq1$;
substitution of $V_j=gf^j$ proves both using
$\sum_{i\geq0}\binom{i+2}{2}f^i=(1-f)^{-3}$.
The recurrences determine each coefficient uniquely. The upper-left
entries therefore have generating function $C$ and the top-row sums
have generating function $g/(1-f)=C^2$.

\begin{corollary}[The floor $-2$ zero-free sibling]\label{cor:thirdpower}
Let $v_N=\mathrm{A006632}(N)$ for $N\geq1$ and let
$V(t)=\sum_{N\geq1}v_Nt^N$. Then
\[
 v_N=b_3(N-1),\qquad V=tC^3,\qquad
 V=\frac{t}{(1-V)^3}.
\]
For $N\geq2$,
\[
 3N(3N-1)(3N-2)v_N
 =8(4N-5)(4N-3)(2N-1)v_{N-1}.
\]
Writing $B=V/t=C^3$, one has $B(0)=1$ and
$[t^n]B^{-n}=-3$ for all $n\geq1$; these conditions determine $B$ uniquely.
\end{corollary}
\begin{proof}
The defining binomial formula in A006632 \cite{oeis7} equals
\eqref{eq:b} at $r=3$, $k=N-1$. Hence $V=tC^3=(C-1)/C$,
which gives its equation. The factorial ratio in Section 3, with
$k=N-2$, gives the recurrence. Lagrange inversion gives
$[t^n]C^{-3n}=-3\binom{n-1}{n-1}=-3$.
The coefficient of $t^n$ in $B^{-n}$ is $-n[t^n]B$ plus a
polynomial in the earlier coefficients, proving uniqueness.
\end{proof}
The three recorded formulas are credited to van Hoeij (11 November 2011;
corrected by Manyama, 28 November 2025), Mathar (7 May 2021), and
Bala (24 July 2025), respectively. They are proved as recorded, rather
than presented as new formulas. Other interpretations and identifications
in A006632 are not covered separately.

\section{Source coverage and formal statements}\label{sec:coverage}
The table covers the posing entries and their first cross-referenced
companions \cite{oeis4,oeis5,oeis6}. Other combinatorial interpretations are distinguished from
identities settled by the present generating functions.
\begingroup\small
\begin{longtable}{@{}p{.32\linewidth}p{.62\linewidth}@{}}
\toprule Source item and locator & Mathematical coverage\\\midrule\endhead
A379464 rev. 20, second Formula & Proved: Mathar recurrence, $n\geq7$.\\
A379463 rev. 16, second Formula & Proved: Mathar recurrence, $n\geq5$.\\
Both entries, Scheuerle programs, 7 Jan. 2025 & Reproved; prior proof Krattenthaler, 2015, Theorem 10.4.5; Prodinger, 2019/2020, Theorem 1 (crediting Selkirk, 2019); Asinowski--Hackl--Selkirk, 2020/2022, Proposition 3; Drube, 2022, Corollary 2.4. Horizontal insertion gives the recorded sums.\\
Both entries, first Formula & Proved: Kotesovec's recorded asymptotics, Corollary~\ref{cor:asymptotic}.\\
A127902 rev. 38, first Formula & Refuted-and-corrected: the printed offset needs $N-1$; $N=4$ is a counterexample. Prior proof of the zero-based sum: Barry, 2015/2016, Example 16; Corollary~\ref{cor:offset}.\\
A127902, Kotesovec's 20 Aug. 2013 formulas & Proved: five-term recurrence for $N\geq5$ and the recorded asymptotic.\\
A127902, definition; Gutkovskiy 1 July 2020 & Proved: series reversion and $Y=x(1+Y+Y^4)$.\\
A127902, Devlin 4 Mar. 2012; Heinz 14 Apr. 2014 & Refuted-and-corrected: distinct vertex labeling fails at two vertices; ordered unlabeled rooted trees work. Proved: Heinz's path interpretation.\\
A069271 rev. 178, Selkirk's 31 Mar. 2020 comments & Reproved; prior proof Prodinger, 2019/2020, Theorem 1: floor $-1$ paths and pairs of full $4$-ary trees, generated by $C^2$.\\
A069271, Mathar 7 June 2013; Eldar 14 Sept. 2025 & Proved: zero-free recurrence and asymptotic, Corollary~\ref{cor:zerofree}.\\
A069271, Bala 10 Sept. 2025 & Refuted-and-corrected: inverse-power value $4(-1)^n$; $n=2$ gives $4$, not $-4$. The divided power formula excludes $m=-2$.\\
A069271, Somos PARI 28 Mar. 2012 & Refuted-and-corrected: remove the extra factor $2$; already false at $n=0$.\\
A069271, remaining algebraic formulas & Proved: square-root, reversion, hypergeometric, logarithmic-derivative and 2026 substitution identities; production matrix, Section~\ref{sec:siblings}. Identifications with further sequence IDs are not covered separately.\\
A002293, Fuss--Catalan formula and powers & Reproved; prior proof Prodinger, 2019/2020, Section 1: $[t^k]C^r=\frac{r}{4k+r}\binom{4k+r}{k}$. Its quartic equation, coefficient convolution, hypergeometric series and asymptotic follow as above.\\
A002293, restricted-word conjecture and other interpretations & Not covered: pattern-avoiding words, maps, forests, tilings, monoids, moment densities and external transform identifications.\\
A069271, specialized plane-tree interpretations & Not covered: the separately defined plane-tree classes in its linked papers.\\
A006632 rev. 152, corrected g.f.; Mathar 7 May 2021; Bala 24 July 2025 & Proved: $V=tC^3=t/(1-V)^3$, the recorded recurrence for $N\geq2$, and the zero-offset inverse-power characterization, Corollary~\ref{cor:thirdpower}. Prior enumeration: the $r=3$ Raney formula. Other interpretations are not covered separately.\\
A240904 rev. 16 & Not covered: it has increments $\{1,3,0,-1,-3\}$ and an upper barrier, unlike the present class.\\
All fixed lower floors & Reproved; prior proof Prodinger, 2019/2020, Section 3, then Drube, 2022, equation (4). Corollary~\ref{cor:allfloors} supplies horizontal insertion.\\
\bottomrule
\end{longtable}\endgroup

The Lean declaration
\href{https://github.com/the-omega-institute/trureturing/blob/03150b25ffad5ced1bebdca2efc8c05b484c1eca/D5/S1/Recurrence/MatharA379464/Result.lean}{\texttt{MatharA379464.Result.result}}
has as its statement the eight-term recurrence for all $n\geq7$, with
the sequence defined as the cardinality of length-$n$ words with increments
$1,0,-3$, total sum zero, and every prefix sum at least $-2$.
In the pinned
\href{https://github.com/the-omega-institute/trureturing/blob/03150b25ffad5ced1bebdca2efc8c05b484c1eca/D5/S1/Words/LatticePaths/MatharA379464Counting.lean}{\texttt{MatharA379464Counting.lean}},
\texttt{three\_positive\_rotations} is Lemma~\ref{lem:cycle} for $r=3$,
and \texttt{binomial\_counting} is \eqref{eq:count} for $r=3$.
Its displayed sum ranges over $0\leq k\leq n$; terms with $4k>n$ vanish
because $\binom{n}{4k}=0$. The other floors, their recurrences, the
all-floor correction and analytic asymptotics have the written proofs given
here. The accompanying checker compares
literal path counts with \eqref{eq:count} for $0\leq n\leq240$ and checks
the polynomial differential identities in exact arithmetic. It also compares
\eqref{eq:allfloors} with literal counts for $0\leq h\leq8$ and
$0\leq n\leq120$, verifies \eqref{eq:powers} for $-5\leq m\leq5$,
$1\leq n\leq16$, and checks the production matrix through power $30$.
These finite comparisons are separate from the all-index written proofs.

\appendix
\section{Recurrence coefficients}\label{sec:coefficients}
In each row below the entry with index $j$ multiplies $a_r(n-j)$.
The initial values and the range $n\geq d_r$ are as in Theorem~\ref{thm:main}.
For the eight-term recurrence ($r=3$):
\begin{align*}
p_{3,0}&=3(n+4)(3n+4)(3n+8),\\
p_{3,1}&=3(-63n^3-297n^2-349n-60),\\
p_{3,2}&=3(189n^3+270n^2-229n-140),\\
p_{3,3}&=15(-63n^3+117n^2+44n-64),\\
p_{3,4}&=689n^3-5372n^2+6946n-1288,\\
p_{3,5}&=(n-4)(201n^2+2767n-3011),\\
p_{3,6}&=-(n-5)(579n+257)(n-4),\\
p_{3,7}&=229(n-5)(n-6)(n-4).
\end{align*}
For the six-term recurrence ($r=2$):
\begin{align*}
p_{2,0}&=3n(3n+4)(n-3)(3n+8),\\
p_{2,1}&=3(-45n^4+54n^3+192n^2-27n-20),\\
p_{2,2}&=9(n-1)(30n^3-72n^2-7n+20),\\
p_{2,3}&=-3(n-1)(n-2)(90n^2-234n+95),\\
p_{2,4}&=-(n-1)(n-2)(n-3)(121n+499),\\
p_{2,5}&=229(n-1)(n-2)(n-3)(n-4).
\end{align*}
For the five-term recurrence ($r=1$, with the offset already changed):
\begin{align*}
p_{1,0}&=3n(3n-4)(3n+4),\\
p_{1,1}&=-3(2n-1)(18n^2-18n-7),\\
p_{1,2}&=3(n-1)(54n^2-108n+47),\\
p_{1,3}&=-54(n-2)(n-1)(2n-3),\\
p_{1,4}&=-229(n-3)(n-2)(n-1).
\end{align*}

\begin{thebibliography}{99}\small
\bibitem{oeis3} The OEIS Foundation Inc.,
\emph{A379464: Paths with steps $(1,1),(1,0),(1,-3)$ staying on or above $y=-2$},
entry by Emely Hanna Li Lobnig, 23 December 2024;
formula contributions by Thomas Scheuerle (7 January 2025),
Vaclav Kotesovec (15 January 2025), and R. J. Mathar (29 January 2025).
Revision 20. \url{https://oeis.org/A379464}.
\bibitem{oeis2} The OEIS Foundation Inc.,
\emph{A379463: Paths with steps $(1,1),(1,0),(1,-3)$ staying on or above $y=-1$},
entry by Emely Hanna Li Lobnig, 23 December 2024;
formula contributions by Thomas Scheuerle (7 January 2025),
Vaclav Kotesovec (15 January 2025), and R. J. Mathar (29 January 2025).
Revision 16. \url{https://oeis.org/A379463}.
\bibitem{oeis1} The OEIS Foundation Inc.,
\emph{A127902: Series reversion of $x/(1+x+x^4)$},
entry by Paul Barry, 5 February 2007; recurrence and asymptotic
by Vaclav Kotesovec, 20 August 2013. Revision 38, 30 May 2026.
\url{https://oeis.org/A127902}.
\bibitem{ac} Philippe Flajolet and Robert Sedgewick,
\emph{Analytic Combinatorics}, Cambridge University Press, 2009.
\href{https://doi.org/10.1017/CBO9780511801655}{doi:10.1017/CBO9780511801655}.
Author edition: \url{https://algo.inria.fr/flajolet/Publications/book.pdf}.
\bibitem{oeis4} The OEIS Foundation Inc.,
\emph{A069271: $\binom{4n+1}{n}2/(3n+2)$}, entry by Henry Bottomley,
12 March 2002. Revision 178, 9 September 2026.
\url{https://oeis.org/A069271}.
\bibitem{oeis5} The OEIS Foundation Inc.,
\emph{A002293: Number of dissections of a polygon}, Fuss--Catalan sequence.
\url{https://oeis.org/A002293}.
\bibitem{oeis6} The OEIS Foundation Inc.,
\emph{A240904: Lattice paths with five allowed increments and two barriers},
entry by Alois P.\ Heinz. Revision 16, 20 December 2020.
\url{https://oeis.org/A240904}.
\bibitem{prodinger} Helmut Prodinger,
\emph{On $k$-Dyck paths with a negative boundary},
arXiv:1912.06930, first version 2019; version 3, 18 August 2020.
Theorem 1 and Section 3, pp.\ 3--8.
\url{https://arxiv.org/abs/1912.06930v3}.
Published in \emph{Journal of Combinatorial Mathematics and Combinatorial
Computing} \textbf{119} (2024), 3--12,
\href{https://doi.org/10.61091/jcmcc119-01}{doi:10.61091/jcmcc119-01}.
\bibitem{drube} Paul Drube,
\emph{Raised $k$-Dyck Paths}, arXiv:2206.01194v1, 2 June 2022.
Theorem 2.3, Corollary 2.4 and equation (4), pp.\ 3--4.
\url{https://arxiv.org/abs/2206.01194v1}.
\bibitem{krattenthaler} Christian Krattenthaler,
\emph{Lattice Path Enumeration}, in M.\ B\'ona (ed.),
\emph{Handbook of Enumerative Combinatorics}, CRC Press, 2015,
pp.\ 589--678. Theorems 10.4.5 and 10.4.7.
\url{https://arxiv.org/abs/1503.05930v3}.
\bibitem{ahs} Andrei Asinowski, Benjamin Hackl and Sarah J.\ Selkirk,
\emph{Down-step statistics in generalized Dyck paths},
\emph{Discrete Mathematics \& Theoretical Computer Science}
\textbf{24}(1) (2022), article 17. Preprint first released 30 July 2020;
version 6, 13 May 2022, Proposition 3 and equation (2).
\url{https://arxiv.org/abs/2007.15562v6}.
\href{https://doi.org/10.46298/dmtcs.7163}{doi:10.46298/dmtcs.7163}.
\bibitem{dejager} Isaac DeJager, Frank Seidl, Madeleine Naquin and Paul Drube,
\emph{Colored Motzkin Paths of Higher Order},
\emph{Journal of Integer Sequences} \textbf{24} (2021), article 21.4.6,
Corollary 11, p.\ 12.
\url{https://cs.uwaterloo.ca/journals/JIS/VOL24/Drube/drube7.pdf}.
\bibitem{barry} Paul Barry,
\emph{Riordan arrays, generalized Narayana triangles, and series reversion},
\emph{Linear Algebra and its Applications} \textbf{491} (2016), 343--385.
Available online 17 November 2015. Example 16, pp.\ 361--362.
\href{https://doi.org/10.1016/j.laa.2015.10.032}{doi:10.1016/j.laa.2015.10.032}.
\bibitem{bf} Cyril Banderier and Philippe Flajolet,
\emph{Basic analytic combinatorics of directed lattice paths},
\emph{Theoretical Computer Science} \textbf{281} (2002), 37--80.
Examples 3 and 6, pp.\ 52 and 64, equations (24), (44)--(45).
\href{https://doi.org/10.1016/S0304-3975(02)00007-5}{doi:10.1016/S0304-3975(02)00007-5}.
\bibitem{oeis7} The OEIS Foundation Inc.,
\emph{A006632: $3\binom{4n-1}{n-1}/(4n-1)$},
revision 152, 22 October 2026.
\url{https://oeis.org/A006632}.
\bibitem{mathar2021} Richard J.\ Mathar,
\emph{P-finite Recurrences From Generating Functions with Roots of Polynomials},
arXiv:2109.02112v1, 5 September 2021, Sections 1--2.
\url{https://arxiv.org/abs/2109.02112v1}.
\end{thebibliography}
\end{document}
